\chapter{2006年北京卷（理）}
\section{单选题}
\begin{problem}
在复平面内，复数 \(\frac{1 + \i}{\i}\) 对应的点位于（\quad）

\choices
    {第一象限}
    {第二象限}
    {第三象限}
    {第四象限}
\end{problem}

\begin{problem}
若 \(\overrightarrow{\alpha}\) 与 \(\overrightarrow{b} - \overrightarrow{c}\) 都是非零向量，则 \(\overrightarrow{a} \cdot \overrightarrow{b} = \overrightarrow{a} \cdot \overrightarrow{c}\) 是 \(\overrightarrow{a} \cdot \overrightarrow{c} \perp (\overrightarrow{b} - \overrightarrow{c})\) 的（\quad）

\choices
    {充分而不必要条件}
    {必要而不充分条件}
    {充分必要条件}
    {既不充分也不必要条件}
\end{problem}

\begin{problem}
在 1, 2, 3, 4, 5 这五个数字组成的没有重复数字的三位数中，各位数字之和为奇数的共有（\quad）

\choices
    {36 个}
    {24 个}
    {18 个}
    {6 个}
\end{problem}

\begin{problem}
平面 \(\alpha\) 的斜线 \(AB\) 交 \(\alpha\) 于点 \(B\)，过定点 \(A\) 的动直线 \(l\) 与 \(AB\) 垂直，且交 \(\alpha\) 于点 \(C\)，则动点 \(C\) 的轨迹是（\quad）

\choices
    {一条直线}
    {一个圆}
    {一个椭圆}
    {双曲线的一支}
\end{problem}

\begin{problem}
已知 \( f(x) = \begin{cases} (3a - 1)x + 4a, & z < 1, \\ 9\lg_{a}x, & z \geqslant 1, \end{cases} \) 是 \((- \infty, +\infty)\) 上的减函数，那么 \( a \) 的取值范围是（\quad）

\choices
    {(0,1)}
    {\(\left(0, \frac{1}{3}\right)\)}
    {\(\left[\frac{1}{7}, \frac{1}{3}\right)\)}
    {\(\left[\frac{1}{2}, 1\right)\)}
\end{problem}

\begin{problem}
在下列四个函数中，满足性质 “对于区间 \( (1,2) \) 上的任意 \( x_{1}, x_{2} (x_{1} \neq x_{2}) \) ，\( |f(x_{1}) - f(x_{2})| < |x_{2} - x_{1}| \) 恒成立” 的只有（\quad）

\choices
    {\(f(x) = \frac{1}{2}\)}
    {\(f(x) = |x|\)}
    {\(f(x) = 2^{x}\)}
    {\(f(x) = x^{2}\)}
\end{problem}

\begin{problem}
设 \( f(0) = 2 + 2^4 + 2^7 + 2^{10} + \cdots + 2^{2n+10} (n \in \mathbf{N}) \)，则 \( f(x) \) 等于（\quad）

\choices
    {\(\frac{2}{5} (8^n - 1)\)}
    {\(\frac{2}{5} (8^{n + 5} - 1)\)}
    {\(\frac{2}{5} (8^{n + 3} - 1)\)}
    {\(\frac{2}{5} (8^{n + 4} - 1)\)}
\end{problem}

\begin{problem}
如图为某三路口交通环岛的简化模型，在某高峰时段，单位时间进山路口 A, B, C，的机动车辆数如图所示，图中 \( x_{1}, x_{2}, x_{3} \) 分别表示该时段单位时间通过路段 \( \overline{AB}, \overline{BC}, \overline{CA} \) 的机动车辆数（假设：单位时间内，在上述路段中，同一路段上驶入与驶出的车辆数相等），则 \( x_{1}, x_{2}, x_{3} \) 的大小关系是（\quad）

\choices
    {\(x_{1} > x_{2} > x_{3}\)}
    {\(x_{1} > x_{3} > x_{2}\)}
    {\(x_{2} > x_{3} > x_{1}\)}
    {\(x_{3} > x_{2} > x_{1}\)}
\end{problem}

\section{填空题}
\begin{problem}
\(\lim_{x\to 1}\frac{x^2 + 3x + 2}{x^2 - 1}\) 的值等于
\end{problem}

\begin{problem}
在 \(\left(\sqrt{x} = \frac{2}{x}\right)^7\) 的展开式中，\(x^2\) 的系数为 \fillinblank{}. (用数字作答)
\end{problem}

\begin{problem}
若三点 \(A(2,2), B(a,0), C(0,b) (ab \neq 0)\) 共线，则 \(\frac{1}{a} + \frac{1}{b}\) 的值等于 \fillinblank{} \fillinblank{}.
\end{problem}

\begin{problem}
在 \(\triangle ABC\) 中，若 \(\sin A: \sin B: \sin C = 5:7:8\)，则 \(\angle B\) 的大小是 \fillinblank{} \fillinblank{}.
\end{problem}

\begin{problem}
已知点 \(P(x, y)\) 的坐标满足条件 \(\left\{ \begin{array}{l} x + y \leqslant 4, \\ y \geqslant x, \\ x \geqslant 1, \end{array} \right.\) 点 \(O\) 为坐标原点，那么 \(|PO|\) 的最小值等于 \fillinblank{}. 最大值等于 \fillinblank{} \fillinblank{}.
\end{problem}

\begin{problem}
已知 \(A, B, C\) 三点在球心为 \(O\)，半径为 \(R\) 的球面上，\(AC \perp BC\)，且 \(AB = R\). 那么 \(A, B\) 两点的球面距离为 ，球心到平面 \(ABC\) 的距离为 \fillinblank{} \fillinblank{}.
\end{problem}

\section{解答题}
\begin{problem}
已知函数 \( f(x) = \frac{1 - \sqrt{2}\sin(2x - \frac{\pi}{3})}{\cos x} \). (1) 求 \( f(x) \) 的定义域. (2) 设 \(\alpha\) 是第四象限的角，且 \(\tan \alpha = -\frac{4}{3}\)，求 \(f(\alpha)\) 的值.
\end{problem}

\begin{problem}
已知函数 \( f(x) = ax^3 + bx^2 + cx \) 在点 \( x_0 \) 处取得极大值 5，其导函数 \( y = f(x) \) 的图象经过点 (1,0)，(2,0). 如图所示，求： (1) \(x_0\) 的值; (2) \(a, b, c\) 的值.
\end{problem}

\begin{problem}
如图，在底面为平行四边形的四棱锥 \(P - ABCD\) 中，\(AB \perp AC\)，\(PA \perp\) 平面 \(ABCD\)，且 \(PA = AB\)，点 \(E\) 是 \(PD\) 的中点. (1) 求证: \(AC \perp PB\); (2) 求证: \(PB //\) 平面 \(AEC\); (3) 求二面角 \(E - AC - B\) 的大小.
\end{problem}

\begin{problem}
某公司招聘员工，指定三门考试课程，有两种考试方案. 方案一：考试三门课程，至少有两门及格为考试通过； 方案二：在三门课程中，随机选取两门. 这两门都及格为考试通过. 假设某应聘者对三门指定课程考试及格的概率分别是 n, b, c，且三门课程考试是否及格相互之间没有影响. (1) 分别求该应聘者用方案一和方案二时考试通过的概率; (2) 试比较该应聘者在上述两种方案下考试通过的概率的大小. (说明理由)
\end{problem}

\begin{problem}
已知点 \(M(-2,0), N(2,0)\)，动点 \(P\) 满足条件 \(|PM| - |PN| = 2\sqrt{2}\). 记动点 \(P\) 的轨迹为 \(W\). (1) 求 \(W\) 的方程; (2) 若 \(A, B\) 是 \(W\) 上的不同两点，\(O\) 是坐标原点，求 \(\overrightarrow{OA} \cdot \overrightarrow{OB}\) 的最小值.
\end{problem}

\begin{problem}
在数列 \(\{a_{n}\}\) 中，若 \(a_{1}, a_{2}\) 是正整数，且 \(a_{n} = |a_{n-1} - a_{n-2}|, n = 3, 4, 5, \cdots\)，则称 \(\{a_{n}\}\) 为“绝对差数列”. (1) 举出一个前五项不为零的“绝对差数列”(只要求写出前十项); (2) 若“绝对差数列” \( \{a_{n}\} \) 中， \( a_{20}=3, a_{21}=0 \) ，数列 \( \{b_{n}\} \) 满足 \( b_{n}=a_{n}+a_{n+1}+a_{n+2}, n=1,2,3,\cdots \) ，分别判断当 \( n\to\infty \) 时， \( a_{n} \) 与 \( b_{n} \) 的极限是否存在. 如果在 \( b_{n} \) 求出其极限值； (3) 证明：任何“绝对差数列”中总含有无穷多个为零的项.
\end{problem}

